%% Chapitre 08 — Coûts, budget et efficacité
%% RÈGLE : uniquement des commandes sémantiques bookforge. Aucun \chapter,
%% \textbf, \emph, \begin{lstlisting}, \vspace... brut.

\chapitre{Coûts, budget et efficacité}

\introChapitre{
  Déléguer a un prix, en jetons comme en temps. Ce chapitre, en partie daté car
  les tarifs bougent, t'apprend à comprendre d'où vient la facture, à estimer le
  coût d'une tâche, et à obtenir plus avec moins. L'efficacité n'est pas
  radiner : c'est viser le bon niveau d'effort pour chaque tâche.
}

\begin{avertissement}
Les prix exacts, les noms de modèles et les paliers de facturation changent
vite. Ce chapitre donne volontairement peu de chiffres : il vise les principes,
qui eux durent. Avant toute décision budgétaire, vérifie la tarification courante
sur la page officielle d'Anthropic et dans ton tableau de bord de consommation.
Ce que tu lis ici t'apprend à raisonner, pas à mémoriser un tarif.
\end{avertissement}

\section{D'où vient le coût}

\paragraphe{Tu as appris au chapitre précédent à vérifier ce que produit l'agent
plutôt qu'à lui faire une confiance aveugle. Cette vérification a un coût, et
l'agent lui-même en a un. Comprendre la facture commence par une unité : le
\terme{jeton}. Un jeton est un morceau de texte, un mot court vaut à peu près un
jeton. Tout ce que l'agent lit et écrit se compte ainsi.}

\paragraphe{On facture deux flux séparés. Les \emphase{jetons d'entrée} sont tout
ce que tu envoies au modèle : ta demande, mais surtout le contexte qui
l'accompagne, fichiers lus, historique de la session, instructions de la mémoire
projet. Les \emphase{jetons de sortie} sont ce que l'agent produit : ses
explications et son code. La sortie coûte en général nettement plus cher par
jeton que l'entrée.}

\paragraphe{Le piège, c'est l'entrée invisible. À chaque tour de boucle, l'agent
ne se contente pas de lire ta nouvelle phrase : il relit tout le contexte
accumulé pour répondre. Une longue session renvoie donc le même historique
encore et encore. C'est le \emphase{contexte rechargé} : trois échanges anodins
dans une conversation de cinquante mille jetons coûtent bien plus que les mêmes
trois échanges dans une session fraîche.}

\begin{note}
La mise en cache du contexte change la donne sur ce point : relire un bloc déjà
vu coûte une fraction du plein tarif. Tu n'as rien à activer, mais ça récompense
un réflexe simple : garder stable le début d'une session plutôt que de remanier
sans cesse les premières instructions.
\end{note}

\section{Estimer avant de lancer}

\paragraphe{Avant de lancer une tâche lourde, prends dix secondes pour estimer
son ordre de grandeur. Pas un chiffre précis : une catégorie. Pose-toi trois
questions. Combien de fichiers l'agent va-t-il devoir lire pour cette tâche ?
Combien de code va-t-il produire ? Combien d'allers-retours, donc combien de
rechargements de contexte, faudra-t-il ?}

\paragraphe{Ces trois réponses te placent dans l'un de trois paniers. Une
question ciblée sur un fichier est \emphase{légère}. Une refonte qui touche
plusieurs modules avec relectures et tests est \emphase{lourde}. Le reste est
\emphase{moyen}. Tu n'as pas besoin de plus de finesse : l'estimation sert à
décider du cadre, pas à tenir une comptabilité.}

\begin{astuce}
Avant une grosse tâche, demande un plan d'abord, exécution ensuite : « décris ton
plan et les fichiers que tu comptes lire, sans rien modifier ». Tu vois le
périmètre réel pour quelques jetons, et tu corriges le tir avant de payer la
boucle complète. Un plan à dix mille jetons qui t'évite une exécution à cent
mille est la meilleure affaire du chapitre.
\end{astuce}

\section{Les gaspillages classiques}

\paragraphe{La plupart des factures salées ne viennent pas de tâches difficiles,
mais de gaspillages évitables. Le premier est la \emphase{session-fleuve} : tu
enchaînes dix sujets sans rapport dans une même conversation, et chaque message
traîne le contexte de tous les précédents. Démarrer une session neuve par sujet
coûte presque toujours moins cher.}

\paragraphe{Le deuxième est le \emphase{contexte gavé} : tu colles un fichier
entier de trois mille lignes quand l'agent n'avait besoin que d'une fonction.
Vise juste, donne le strict nécessaire. Tu as vu au chapitre sur le contexte que
trop d'information dégrade aussi la qualité ; ici, elle dégrade en plus le
portefeuille.}

\paragraphe{Le troisième est la \emphase{boucle qui s'acharne} : l'agent ne
trouve pas, réessaie, relit, repart, et la facture grimpe sans que le problème
avance. Quand tu sens ce surplace, coupe. Reformule la tâche, donne l'indice qui
manque, ou reprends la main. Laisser tourner par paresse est le gaspillage le
plus sournois.}

\begin{liste}
  \element{Une session par sujet : on n'empile pas des thèmes sans rapport.}
  \element{Le strict nécessaire en entrée : un extrait précis vaut mieux qu'un
  fichier entier.}
  \element{On coupe une boucle qui patine au lieu de la laisser brûler des jetons.}
\end{liste}

\section{Le bon niveau d'effort}

\paragraphe{L'efficacité n'est pas de toujours dépenser le moins possible. C'est
d'aligner l'effort sur l'enjeu. Renommer une variable et refondre une couche
d'authentification ne méritent ni le même soin ni le même budget. Sous-investir
sur une tâche critique coûte cher en bugs ; sur-investir sur une broutille
gaspille sans rien gagner.}

\paragraphe{Deux leviers règlent cet effort. Le premier est le \emphase{choix du
modèle} : un modèle plus puissant raisonne mieux mais coûte plus par jeton.
Réserve-le aux tâches où le raisonnement compte vraiment, et confie le travail
mécanique et répétitif à un modèle plus léger. Le second est l'\emphase{ampleur
de la consigne} : demander un plan détaillé, des tests et une justification a du
sens pour un changement risqué, beaucoup moins pour une correction triviale.}

\paragraphe{La règle tient en une phrase : \important{dépense à la hauteur du coût
d'une erreur.} Là où se tromper se paie cher, mets le paquet, modèle fort,
vérification serrée. Là où une erreur se corrige en dix secondes, vas-y au plus
simple et au moins cher.}

\section{Suivre sa consommation}

\paragraphe{On ne pilote bien que ce qu'on mesure. Selon ta formule, ta
consommation se suit soit dans un tableau de bord d'usage, soit via les
indicateurs affichés en cours de session. Prends l'habitude d'y jeter un œil :
non pour culpabiliser, mais pour calibrer ton intuition. Au bout de quelques
jours, tu sauras d'instinct ce que coûte une tâche moyenne.}

\paragraphe{Fixe-toi ensuite des repères. Sur une formule à l'usage, une limite
de dépense mensuelle t'évite la mauvaise surprise et te force à prioriser. Sur un
forfait, c'est ton quota de session qu'il faut apprendre à gérer pour ne pas
buter sur le plafond au pire moment.}

\begin{avertissement}
Méfie-toi des tâches lancées et oubliées en arrière-plan. Une boucle qui tourne
sans surveillance peut consommer beaucoup avant que tu t'en rendes compte. Pose
un garde-fou : un périmètre borné, une limite, et un coup d'œil régulier. L'agent
n'a aucune notion de ce qu'il dépense ; cette vigilance, c'est ton rôle.
\end{avertissement}

\section{Humain ou agent ?}

\paragraphe{Dernière question, la plus saine : cette tâche, l'agent doit-il
vraiment la faire ? Déléguer n'est rentable que si le temps gagné dépasse le coût
en jetons plus le temps de vérification. Parfois, le calcul penche pour toi.}

\paragraphe{L'agent gagne quand la tâche est volumineuse mais mécanique : migrer
cinquante fichiers vers une nouvelle interface, écrire une batterie de tests
répétitifs, traduire de la documentation. Beaucoup de sortie, peu de jugement, et
un résultat facile à vérifier en lot : c'est son terrain.}

\paragraphe{Tu gagnes, toi, quand la tâche est minuscule et que tu sais déjà quoi
faire. Corriger un caractère, ajuster une valeur de configuration, renommer un
fichier : le temps d'écrire la consigne et de relire le \emphase{diff}, tu
l'aurais fait deux fois à la main. Le bon réflexe n'est pas « tout déléguer »,
mais « déléguer ce qui le mérite ».}

\resumeChapitre{
  \element{\important{Jetons} : tu paies l'entrée (contexte rechargé à chaque tour) et la sortie (code produit) — la sortie est plus chère par jeton, mais l'entrée est le piège invisible des longues sessions.}
  \element{\important{Estimer avant de lancer} : demande d'abord un plan sans exécution — quelques jetons pour voir le périmètre réel avant de payer la boucle complète.}
  \element{\important{Trois gaspillages} : la session-fleuve sur des sujets sans rapport, le contexte gavé de fichiers inutiles, et la boucle qui s'acharne sans avancer — coupe-les sans hésiter.}
  \element{\important{Aligner l'effort sur l'enjeu} : modèle fort et vérification serrée pour les tâches risquées, modèle léger et consigne simple pour les tâches mécaniques.}
  \element{\important{Mesurer} : suis ta consommation pour calibrer ton intuition ; fixe une limite mensuelle pour prioriser et éviter la mauvaise surprise.}
}

\paragraphe{Maîtriser le coût, c'est gagner le droit de répéter ses meilleurs
gestes sans hésiter. Une fois que tu sais ce que vaut une tâche et comment la
cadrer, l'étape suivante est naturelle : transformer les gestes qui reviennent en
raccourcis fiables. C'est l'objet du chapitre suivant, où l'on industrialise
commandes et automatisations.}
